\problem{}
An array $A$ of length $n$ contains a random permutation of the distinct integers $1, 2, \dots, n$. All $n!$ permutations are equally likely.

An index $i$ is called a local maximum if the element $A[i]$ is greater than its neighbors. Specifically: \begin{itemize} \item For $1 < i < n$, index $i$ is a local maximum if $A[i-1] < A[i]$ and $A[i] > A[i+1]$. \item For the boundaries, index $1$ is a local maximum if $A[1] > A[2]$, and index $n$ is a local maximum if $A[n] > A[n-1]$. \end{itemize}

Let $M$ be the random variable representing the total number of local maxima in the array. Determine the expectation $\mathbb{E}[M]$.

\noindent \textbf{Solution:}



Define indicator random variables for each index \(i\):
\[
X_i = \begin{cases}
1, & \text{if position } i \text{ is a local maximum},\\
0, & \text{otherwise}.
\end{cases}
\]
Then \(M = \sum_{i=1}^n X_i\), and by linearity of expectation,
\[
\mathbb{E}[M] = \sum_{i=1}^n \mathbb{E}[X_i] = \sum_{i=1}^n \Pr(X_i = 1).
\]

We compute the probability that position \(i\) is a local maximum.

\noindent\textbf{Boundary positions:}
\begin{itemize}
    \item For \(i = 1\): We require \(A[1] > A[2]\). Since the permutation is uniformly random, \(A[1]\) and \(A[2]\) are two distinct random numbers, and all orderings are equally likely. Hence,
    \[
    \Pr(X_1 = 1) = \frac{1}{2}.
    \]
    \item For \(i = n\): Similarly, \(\Pr(X_n = 1) = \Pr(A[n] > A[n-1]) = \frac{1}{2}\).
\end{itemize}

\noindent\textbf{Internal positions \(2 \leq i \leq n-1\):}
We require \(A[i-1] < A[i]\) and \(A[i] > A[i+1]\). Consider the three consecutive values \(A[i-1], A[i], A[i+1]\). They are three distinct numbers chosen uniformly from \(\{1,\dots,n\}\), and all \(3! = 6\) orderings are equally likely. For \(A[i]\) to be the largest among the three, it must be the maximum. The probability that the middle element is the largest is \(\frac{2}{6} = \frac{1}{3}\). Thus,
\[
\Pr(X_i = 1) = \frac{1}{3}.
\]

Therefore, for \(n \ge 2\),
\[
\mathbb{E}[M] = \frac{1}{2} + \frac{1}{2} + (n-2)\cdot\frac{1}{3} = 1 + \frac{n-2}{3} = \frac{n+1}{3}.
\]

(For \(n = 1\), the single element is trivially a local maximum, so \(\mathbb{E}[M] = 1\), but the formula \(\frac{n+1}{3}\) does not apply.)

Hence, the expected number of local maxima is \({\frac{n+1}{3}}\) for \(n \ge 2\).



\newpage